% Folder 33, batch 02 — archivists' pages 21–40.
% Pass: Opus 5.5 (claude-opus-5-5), 2026-10-03 — first pass, unchecked against the pages by a human.
%
% Only the odd pages (21, 23, …, 39) carry his hand. The even pages are the
% typed rectos of unrelated typescripts (a Bourbaki-style text on root systems
% numbered « n° 391 », and an English typescript on presheaves and stalks) and
% are skipped. His own section numbers (15) to 19) are recorded as headings.
\documentclass[11pt,a4paper]{article}
\input{../preamble/grothendieck}

\folder{33}
\batch{2}
\pages{21}{40}
\foldertitle{SGA 1965. MS [Manuscrit] brouillon : notes manuscrites (s.d.).}
\dating{1965-[vers 1971]}
\watermark{Édition de démonstration}

\begin{document}

\section{15) Classe fondamentale locale}

\page{21}
\note{le numéro « 15 », cerclé, est de sa main, en tête de la page ; l'exposé numéroté qui précède se trouve avant ce lot.}

On fixe un entier $n$ premier aux car. résiduelles.

\emph{Classe fondamentale locale}

Supposons $Y \hookrightarrow X$ « intersection complète de dim. relative $1$ » et supposons d'abord
$Y = V(\varphi)$, $\varphi$ section régulière de $\mathcal{O}_X$. Alors si
$i : X - Y = U \to X$, alors \struck{$i_*(\mathcal{O}_U) \simeq \mathcal{O}_X[1/\varphi]$}, \struck{$i_*(\mathcal{O}_U^*) \to$}
\[
  \varphi \in \Gamma\bigl(i_*(\mathbb{G}_{m,U})\bigr)
\]
\struck{$\to\ \in \Gamma(i_*(\mathbb{G}_{m,U}^*))$}

\struck{Donc} D'ailleurs $\varphi$ est déterminé mod multiplication par un \struck{\ill{}}
$\in \Gamma(X, i_*(\mathbb{G}_{m,X}))$, donc $\varphi$ définit une section \struck{\ill{}} canonique de
\struck{$i_*(\mathcal{O}_U^*)/$} $i_*(\mathbb{G}_{m,U})/\mathbb{G}_{m,X}$. Or on a la suite exacte canonique
\[
  0 \to \underline{H}^0_Y(\mathbb{G}_{m,X}) \to \mathbb{G}_{m,X} \to i_*(\mathbb{G}_{m,U})
  \xrightarrow{\ d\ } \underline{H}^1_Y(\mathbb{G}_{m,X}) \to 0
\]
\note{sous $\underline{H}^0_Y(\mathbb{G}_{m,X})$ il écrit « $= 0$ » verticalement, et le $0$ final est placé sous $\underline{H}^1_Y(\mathbb{G}_{m,X})$, au bout d'une flèche verticale.}

qui donne ici \struck{\ill{}}
\[
  \underline{H}^1_Y(\mathbb{G}_{m,X}) \xrightarrow{\ \sim\ } i_*(\mathbb{G}_{m,U})/\mathbb{G}_{m,X}
\]
et montre que $d\varphi$ est une section canonique de $\underline{H}^1_Y(\mathbb{G}_{m,X})$, que nous
\uncertain{noterons} \struck{\ill{}} $\beta_{Y/X}$. D'autre part, on a la suite exacte de Kummer
\[
  0 \to \mu_{n,X} \to \mathbb{G}_{m,X} \xrightarrow{\ n\ } \mathbb{G}_{m,X} \to 0
\]

\page{23}
d'où
\[
  \partial_n : \underline{H}^1_Y(\mathbb{G}_{m,X}) \longrightarrow \underline{H}^2_Y(\mu_{n,X})
\]
et nous pouvons \struck{\ill{}} poser
\[
  \gamma^{(n)}_{Y/X} = \partial_n(\beta_{Y/X}) \in \Gamma\bigl(X, \underline{H}^2_Y(\mu_{n,X})\bigr)
\]
\note{le $X$ de $\Gamma(X, \dots)$ paraît surchargé ; la lecture $Y$ n'est pas exclue.}

\uncertain{en} laissant tomber les exposants $n$ \ill{} \ill{} \ill{} \ill{} \ill{}
\struck{la définition de} Cette section \struck{satisfait} \ill{} \ill{} : la définition
\uncertain{au} Zariski \ill{} $=$ \ill{} \ill{} \ill{} base \ill{}
$X' \to X$ tel que $Y'$ soit une « int.\ complète » (i.e.\ $\varphi'$ section régulière) et que \ill{} \ill{}

pour \ill{} définir $\gamma^{(n)}_{Y/X}$ en général \add{maintenant} (\uncertain{sans répétition} \ill{}). Si $Y$
\uncertain{intersection} complète de codim.\ $d$, \ill{} localement $Y$ \ill{} \ill{} de la \struck{forme}
\[
  Y = Y_1 \cap \dots \cap Y_d
\]
et on pose
\[
  \gamma_{Y/X} = \gamma_{Y_1/X} \cdot \ldots \cdot \gamma_{Y_d/X}
  \in \Gamma\bigl(Y, \underline{H}^{2d}_Y(\mu_n^{\otimes d})\bigr)
\]
\note{au-dessus du premier point de produit, un signe biffé (peut-être « cup »).}

(cup-produit), et on prouve que cela ne dépend pas du choix des $Y_i$, donc se définit globalement.

Propriétés formelles :

\page{25}
\marginal{propriétés caractéristiques}
\note{la note marginale, écrite en biais, embrasse d'une accolade les propriétés (i) à (iii).}

(i) Compatibilité avec chgt de base \add{admissible}.

(ii) Si $Y = V(\varphi)$, $\gamma^{(n)}_{Y/X} = d(\beta^{(n)}_{Y/X})$ comme \add{dessus}.

(iii) $\gamma^{(n)}_{Y.Y'/X} = \gamma^{(n)}_{Y/X} \cdot \gamma^{(n)}_{Y'/X}$ si $Y$, $Y'$ « sans comp.\ excédentaire ».

(iv) Compatibilité par \uncertain{excision}.

(\textbf{NB} Le « bon » \uncertain{dieu} se trouve dans $\Gamma(Y, \underline{H}^{2d}_Y(\mathbb{Z}_\ell(d)))$, il y a un \ill{} pour tt $\ell$
premier aux car.\ résiduelles de $Y$).

\emph{Application à la construction de classes de cohom.\ associées à \uncertain{certains} cycles.}

Supposons
\[
  \underline{H}^i_Y\bigl((\mathbb{Z}/n\mathbb{Z})_X\bigr) = 0 \quad\text{pour } i < 2d
\]
d'où également
\[
  \underline{H}^i_Y(\mu_{n,X}^{\otimes d}) = 0 \quad\text{pour } i < 2d .
\]
C'est la situation du th.\ de pureté, vrai
\begin{enumerate}
  \item[(i)] $X$, $Y$ lisses sur une base $S$
  \item[(ii)] $Y$, $X$ réguliers, $X$ excellent, car.\ nulle
  \item[(iii)] $d = 1$, $X$ normal.
\end{enumerate}

\page{27}
Alors la suite spectrale implique
\[
  H^i_Y(X, \mu_{n,X}^{\otimes d}) = 0 \quad\text{si } i < 2d
\]
\[
  H^{2d}_Y(X, \mu_{n,X}^{\otimes d}) \simeq \Gamma\bigl(Y, \underline{H}^{2d}_Y(\mu_{n,X}^{\otimes d})\bigr)
\]
\note{dans le dernier membre, un $\mathbb{Z}/n$ biffé précède $\mu_{n,X}^{\otimes d}$.}

d'où un \struck{bon cocycle} \add{élément} canonique
\[
  \gamma^{(n)}_{Y/X} \in H^{2d}_Y(X, \mu_{n,X}^{\otimes d})
\]
[qu'on peut maintenant spécialiser dans $H^{2d}(X, \mu_{n,X}^{\otimes d})$, ou si $X$ est propre sur un $S$
dans $H^{2d}_{\Phi_{Y/X}}(X, \mu_{n,X})$.]
\note{l'indice $\Phi_{Y/X}$ du dernier groupe est glosé par une flèche : « famille de supports » ; il est écrit sous un $H^{2d}$ dont la lecture de l'indice reste incertaine.}

N.B.\ ($\beta_{Y/X} \in H^1_Y(X, \mathbb{G}_{m,X})$ et $\gamma^{(n)}_{Y/X} = \partial_n(\beta_{Y/X})$, en particulier \ill{}
\ill{} $H^2_Y(X, \mathbb{G}_{m,X})$ est nulle).

\emph{Exemple 1.} $X$ lisse sur $k$ \add{sép.\ clos}, \struck{$x$ pt} dim relative $d$, \add{$x$ pt rat.\ de $X$ sur $k$, d'où}
\struck{il résulte des \ill{} \ill{} th.\ de} \ill{}
\[
  \gamma_{x/X} \in H^{2d}_{!}(X, \mu_{n,X}^{\otimes d}) .
\]
Je \uncertain{vais} \uncertain{résulter} de \ill{} le théorème qui \ill{} \ill{} que \struck{\ill{}} cet élément ne dépend pas de $x$
si $X$ est géom.\ connexe, d'où un élément de $H^{2d}_{!}(X, \mu_{n,X}^{\otimes d})$ et un \uncertain{dépend} pas du choix de $x$.
\note{ce passage est souligné par lui depuis « d'où un élément » jusqu'à « choix de $x$ » ; la syntaxe de la fin de phrase n'est pas sûre.}

C'est ce qu'on avait vérifié directement dans le cas où $d = 1$, \add{nous ramenant} au cas de $C$
en \struck{la courbe} \ill{} \ill{} complète où on dispose \struck{déjà}

\page{29}
(par le th.\ de Kummer global) d'un isom.\ canonique
\[
  t : H^2(X, \mu_{n,X}) \xrightarrow{\ \sim\ } \mathbb{Z}/n\mathbb{Z} .
\]
On avait vérifié simplement
\[
  t(\gamma_{x/X}) = 1 .
\]

\emph{Exemple 2.} $X$ et $Y$ lisses sur $S$, $Y$ de codim.\ $d$ dans $X$ en chaque pt. Alors on est dans les conditions
favorables, qui impliquent \add{non} \uncertain{aussi} \ill{}

\begin{tikzcd}[column sep=small]
  \Gamma\bigl(Y, \underline{H}^{2d}_Y(\mu_{n,X}^{\otimes d})\bigr) \arrow[d, "\wr"] & \\
  H^{2d}_Y(X, \mu_{n,X}^{\otimes d}) \arrow[r] & \Gamma\bigl(S, R^{2d}_{!} f_{*}(\mu_{n,X}^{\otimes d})\bigr)
\end{tikzcd}

\note{à la première ligne, une flèche biffée menait de $\Gamma(Y, \dots)$ à $H^0(S, R^{2d}_{!} f_{*}(\mu\dots$, abandonné ; au début du second terme de la seconde ligne, un signe raturé.}

\section{16) Classes fondamentales globales}

\page{31}
On fixe un entier $n > 0$, et on \ill{}
\struck{Soit $f : X \to Y$ morphisme} \add{\uncertain{morphismes}} \ill{} faisceaux de $n$-torsion
(i.e.\ Modules sur $A = \mathbb{Z}/n\mathbb{Z}$), et les \ill{} premiers à \ill{} \ill{} les \uncertain{préschémas} où $n$ est
\uncertain{inversible}, i.e.\ tels que $n \cdot 1_{\mathcal{O}_X}$ est inversible dans $\mathcal{O}_X$.

Si $f : X \to Y$ est un morphisme lisse, de dim.\ relative $d$, ou disons qu'il est \add{purement} \add{de} dim.\ $d$
\add{\uncertain{disons}} dans les fibres \ill{} \ill{} \ill{},
\[
  T^{(n)}_{X/Y} \quad\text{ou simplement}\quad T_{X/Y}
\]
\ill{} confusion possible, le faisceau $\mu_{n,X}^{\otimes d}$. [Si $f$ est lisse de dim.\ \add{relative} \ill{}
quelconque, alors on a \ill{} \ill{} \ill{} \ill{} \ill{} $T_{X/Y}$ \ill{} \ill{} \ill{} \ill{} \ill{} \ill{}.]

Nous regardons par la suite des morphismes lisses \emph{compactifiables} \emph{purement de dim.\ $d$}
$f : X \to Y$. Pour un tel morphisme, on peut définir un \struck{\ill{}} \add{homom.}\ canonique

\page{33}
\[
  \mathrm{Tr}_f : R^{2d}_{!} f_{*}(T_{X/Y}) \longrightarrow A_Y
\]
\note{l'$A$ de $A_Y$ est écrit sur un premier $A$ raturé.}

Tel que les propriétés suivantes :

(i) Exact.\ de base \ill{}

(ii) Transitivité pour un \struck{\ill{}} \add{composé} $gf$
[NB.\ $R^{2(d+d')}_{!}(gf)_{*} = (R^{2d}_{!} g)(R^{2d'}_{!} f)$]
\note{l'exposant du premier membre a été retouché.}

(iii) Cas où $f$ est un \emph{isom.}

(iv) \struck{\ill{}} Cas de $\mathbb{P}^1_S \to S$
[avec $S$ spectre d'un corps \ill{} \ill{} sép.\ clos]
\struck{le \uncertain{diviseur}} l'isom.\ canonique fourni par le théorème de Kummer.

Les conditions \emph{caractérisent} la \uncertain{donnée} $\mathrm{Tr}_f$. De plus

(v) Si $f$ est à fibres connexes non vides, $\mathrm{Tr}_f$ est un isom.

(vi) Compatibilité par la \uncertain{dualité}.

\struck{\ill{}}

\section{17) Relations entre classes fond.\ locales et classes fondamentales globales}

\page{35}
\begin{tikzcd}
  Y \arrow[r, hook, "i", "(d)"'] \arrow[d, "g"', "(r)"] & X \arrow[dl, "f", "(r+d)"'] \\
  S &
\end{tikzcd}

\note{les étiquettes $(d)$, $(r)$, $(r+d)$ notent les codimension et dimensions relatives ; l'étiquette $(r)$ est surchargée.}

\struck{$T_{Y/X} \otimes T_{Y/S} \simeq T_{X/S}|Y$, \quad $T_{Y/S} \simeq \underline{H}^{2d}_Y(T_{X/S})$,
\quad $R^{2r}_{!} f(T_{Y/S}) \simeq$ \ill{}, \quad $T_{Y/S} \simeq$ \ill{}}
\note{ces formules, encadrées, sont barrées de longs traits obliques ; elles sont transcrites telles qu'on les lit.}

\[
  \left\{
  \begin{aligned}
    \underline{H}^i_Y(T_{X/S}) &\simeq 0 \quad\text{si } i \neq 2d\\
    \underline{H}^{2d}_Y(T_{X/S}) &\simeq T_{Y/S}
  \end{aligned}
  \right.
\]
\[
  \Longrightarrow\quad
  \left\{
  \begin{aligned}
    R^{i}_{!} g_{*Y}(T_{X/S}) &= 0 \quad\text{si } i < 2d\\
    R^{2r}_{!} g_{*Y}(T_{X/S}) &\simeq R^{2r} g_{*Y}\bigl(\underline{H}^{2d}_Y(T_{X/S})\bigr)
  \end{aligned}
  \right.
\]
\[
  R^{2r} g_{*Y}\bigl(\underline{H}^{2d}_Y(T_{X/S})\bigr)
  \simeq R^{2r+2d}_{\substack{\text{supports propres}/S\\ \text{dans } Y}} f_{*}(T_{X/S})
\]
\note{dans la première accolade de droite, un exposant raturé avant $i$ ; dans la seconde ligne, des surcharges sur l'exposant $2r$ et, sous le $\simeq$, un signe raturé.}

\begin{tikzcd}[column sep=small, row sep=small, nodes={font=\scriptsize}]
  & R^{2r+2d}_{\text{supp.\ propres}/S,\ \subset Y} f_{*}(T_{X/S}) \arrow[d] \\
  (\mathbb{Z}/n\mathbb{Z})_S & R^{2r+2d}_{!} g_{*X}(T_{X/S}) \arrow[l]
\end{tikzcd}

\note{la flèche oblique de la page part de la ligne $R^{2r+2d}_{!} g_{*X}(T_{X/S})$ vers $(\mathbb{Z}/n\mathbb{Z})_S$ ; le $\mathbb{Z}$ y est écrit sur un premier signe raturé. Les indices $g_{*Y}$, $g_{*X}$ sont lus tels quels.}

Compatibilité avec les \uncertain{deux} traces … S'exprime plus aisément quand on dispose \add{de la}
catégorie \add{dérivée}, et \uncertain{ils} nous \ill{} \uncertain{donnent} une forme plus \add{complète}.

18) \emph{Catégories dérivées} ; application aux $R_{!} f_{*}$.

19) \emph{Th.\ de dualité globale}.

\note{ces deux derniers titres annoncent la suite ; le 18) est repris à la page suivante.}

\section{18) Catégories dérivées}

\page{37}
Les foncteurs entre cat.\ triangulées \struck{\ill{}} \add{commutent} aux foncteurs translation et transforment
triangles distingués en \add{triang.\ dist.}\ \ill{} \uncertain{itou} …

$D(C)$, $D^-(C)$, $D^b(C)$ importants \ill{} sans qu'il y ait assez de projectifs.

\[
  Lf^* : D(Y) \longrightarrow D(X) \quad\text{quand on a } f : X \to Y
\]
\[
  D^+(Y) \to D^+(X) \quad \text{morphismes plats}
\]
\[
  D^-(Y) \to D^-(X) \quad \text{à topos annelés}
\]
\[
  D^+(Y) \to D^+(X) \quad \text{définition directe via résolutions injectives}
\]
\emph{ou} définition indirecte par « \uncertain{passage au quotient} ».

[Pour calculer $Lf^*(K)$, pas besoin de remplacer $K$ par une résolution injective]
\note{lecture de la phrase entre crochets très incertaine ; seuls « calculer », « $Lf^*(K)$ » et « résolution » sont sûrs.}

\textbf{NB.} Si $f$ \uncertain{n'est pas plat}, \ill{} on ne peut \ill{}
\[
  f^* : D^+(Y) \to D^+(X)
\]
(\ill{} \ill{} \ill{} \ill{} \ill{} bonnes propriétés). Mais dans la \uncertain{suite} de \uncertain{Verdier}, \ill{} explique
comment construire un « dérivé gauche » de $f^*$,
\[
  Lf^* : D^-(Y) \to D^-(X),
\]
en utilisant \ill{} \ill{} \ill{} \uncertain{projectives} mais des résolutions plates. On \ill{} que si $f$ est
de « Tor-dim.\ finie » on définit
\[
  D^+(Y) \to D^+(X), \qquad D^b(Y) \to D^b(X) .
\]
\marginal{Rien n'empêche de construire un \uncertain{adjoint à droite}, $\mathbb{R}f^{\times}$, mais il n'a pas la \uncertain{variance intuitive}.}

\page{39}
Définition de
\[
  \mathbb{R}\mathrm{Hom}(F^{\bullet}, G^{\bullet}) = \mathrm{Hom}^*\bigl(F^{\bullet}, C^{\bullet}(G^{\bullet})\bigr) .
\]
\struck{\ill{}} Si \ill{} \ill{} des Modules sur un Topos,
\[
  \mathbb{R}\underline{\mathrm{Hom}}(F^{\bullet}, G^{\bullet})
\]
lié au précédent par
\[
  \boxed{\ \mathbb{R}\mathrm{Hom}(F^{\bullet}, G^{\bullet})
  = \mathbb{R}\Gamma_X\, \mathbb{R}\underline{\mathrm{Hom}}(F^{\bullet}, G^{\bullet})\ }
\]
\[
  \left\{
  \begin{aligned}
    &\mathrm{Hom}(F^{\bullet}, C^{\bullet}) \text{ est aussi noté } \mathrm{Ext}^0(F^{\bullet}, C^{\bullet})\\
    &\mathrm{Hom}(F^{\bullet}, C^{\bullet}(i)) \simeq \mathrm{Hom}(F(-i), C^{\bullet}) = \mathrm{Ext}^i(F^{\bullet}, C^{\bullet})
  \end{aligned}
  \right.
\]
Notation $\underline{\mathrm{Ext}}^i(F^{\bullet}; C^{\bullet})$ faisceaux.
\[
  \mathrm{Ext}^*(F^{\bullet}; C^{\bullet}) \Longleftarrow
  E_2^{pq} = H^p\bigl(X, \underline{\mathrm{Ext}}^q(F^{\bullet}; C^{\bullet})\bigr)
\]
\[
  {}'E_2^{pq} = H^p\bigl(X, \mathrm{Ext}^q(F^*, C^*)\bigr)
\]
\marginal{pris terme à terme}
\note{la note marginale renvoie par deux traits à $F^*$ et $C^*$ ; l'argument se poursuit vraisemblablement au-delà de ce lot.}

\end{document}
